\documentclass[12pt,letterpaper]{article}
\usepackage{graphicx,textcomp}
\usepackage{natbib}
\usepackage{setspace}
\usepackage{fullpage}
\usepackage{color}
\usepackage[reqno]{amsmath}
\usepackage{amsthm}
\usepackage{fancyvrb}
\usepackage{amssymb,enumerate}
\usepackage[all]{xy}
\usepackage{endnotes}
\usepackage{lscape}
\newtheorem{com}{Comment}
\usepackage{float}
\usepackage{hyperref}
\newtheorem{lem} {Lemma}
\newtheorem{prop}{Proposition}
\newtheorem{thm}{Theorem}
\newtheorem{defn}{Definition}
\newtheorem{cor}{Corollary}
\newtheorem{obs}{Observation}
\usepackage[compact]{titlesec}
\usepackage{dcolumn}
\usepackage{tikz}
\usetikzlibrary{arrows}
\usepackage{multirow}
\usepackage{xcolor}
\newcolumntype{.}{D{.}{.}{-1}}
\newcolumntype{d}[1]{D{.}{.}{#1}}
\definecolor{light-gray}{gray}{0.65}
\usepackage{url}
\usepackage{listings}
\usepackage{color}

\definecolor{codegreen}{rgb}{0,0.6,0}
\definecolor{codegray}{rgb}{0.5,0.5,0.5}
\definecolor{codepurple}{rgb}{0.58,0,0.82}
\definecolor{backcolour}{rgb}{0.95,0.95,0.92}

\lstdefinestyle{mystyle}{
	backgroundcolor=\color{backcolour},   
	commentstyle=\color{codegreen},
	keywordstyle=\color{magenta},
	numberstyle=\tiny\color{codegray},
	stringstyle=\color{codepurple},
	basicstyle=\footnotesize,
	breakatwhitespace=false,         
	breaklines=true,                 
	captionpos=b,                    
	keepspaces=true,                 
	numbers=left,                    
	numbersep=5pt,                  
	showspaces=false,                
	showstringspaces=false,
	showtabs=false,                  
	tabsize=2
}
\lstset{style=mystyle}
\newcommand{\Sref}[1]{Section~\ref{#1}}
\newtheorem{hyp}{Hypothesis}

\title{Problem Set 3}
\date{Due: March 25, 2026}
\author{Rosalie LaCroix}


\begin{document}
	\maketitle
	\section*{Instructions}
	\begin{itemize}
	\item Please show your work! You may lose points by simply writing in the answer. If the problem requires you to execute commands in \texttt{R}, please include the code you used to get your answers. Please also include the \texttt{.R} file that contains your code. If you are not sure if work needs to be shown for a particular problem, please ask.
\item Your homework should be submitted electronically on GitHub in \texttt{.pdf} form.
\item This problem set is due before 23:59 on Wednesday March 25, 2026. No late assignments will be accepted.
	\end{itemize}

	\vspace{.25cm}
\section*{Question 1}
\vspace{.25cm}
\noindent We are interested in how governments' management of public resources impacts economic prosperity. Our data come from \href{https://www.researchgate.net/profile/Adam_Przeworski/publication/240357392_Classifying_Political_Regimes/links/0deec532194849aefa000000/Classifying-Political-Regimes.pdf}{Alvarez, Cheibub, Limongi, and Przeworski (1996)} and is labelled \texttt{gdpChange.csv} on GitHub. The dataset covers 135 countries observed between 1950 or the year of independence or the first year forwhich data on economic growth are available ("entry year"), and 1990 or the last year for which data on economic growth are available ("exit year"). The unit of analysis is a particular country during a particular year, for a total $>$ 3,500 observations. 

\begin{itemize}
	\item
	Response variable: 
	\begin{itemize}
		\item \texttt{GDPWdiff}: Difference in GDP between year $t$ and $t-1$. Possible categories include: "positive", "negative", or "no change"
	\end{itemize}
	\item
	Explanatory variables: 
	\begin{itemize}
		\item
		\texttt{REG}: 1=Democracy; 0=Non-Democracy
		\item
		\texttt{OIL}: 1=if the average ratio of fuel exports to total exports in 1984-86 exceeded 50\%; 0= otherwise
	\end{itemize}
	
		\noindent I first prepared the data set. I turned the outcome variable, GDPWdiff from a continuous variable into a discrete variable with levels "positive", "no change", and "negative", corresponding to the sign and value of the numerical difference in GDP between years. I then factored this variable.
	
	\lstinputlisting[language=R, firstline=41, lastline=51]{PS03_RL.R}
	
	
	
\end{itemize}
\newpage
\noindent Please answer the following questions:

\begin{enumerate}
	\item Construct and interpret an unordered multinomial logit with \texttt{GDPWdiff} as the output and "no change" as the reference category, including the estimated cutoff points and coefficients.
	
	\noindent\textbf{My Answer:}
	
	\noindent To run an unordered multinomial logit, I set the reference category for GDPWdiff as "no change".
	
	\lstinputlisting[language=R, firstline=55, lastline=60]{PS03_RL.R}
	
\begin{table}[!htbp] \centering 
	\caption{Table of Coefficients: Unordered Multinomial Logit} 
	\label{} 
	\begin{tabular}{@{\extracolsep{5pt}}lcc} 
		\\[-1.8ex]\hline 
		\hline \\[-1.8ex] 
		& \multicolumn{2}{c}{\textit{Dependent variable:}} \\ 
		\cline{2-3} 
		\\[-1.8ex] & negative & positive \\ 
		\\[-1.8ex] & (1) & (2)\\ 
		\hline \\[-1.8ex] 
		Oil & 4.784 & 4.576 \\ 
		& (6.885) & (6.885) \\ 
		& & \\ 
		Democracy & 1.379$^{*}$ & 1.769$^{**}$ \\ 
		& (0.769) & (0.767) \\ 
		& & \\ 
		Constant & 3.805$^{***}$ & 4.534$^{***}$ \\ 
		& (0.271) & (0.269) \\ 
		& & \\ 
		\hline \\[-1.8ex] 
		Akaike Inf. Crit. & 4,690.770 & 4,690.770 \\ 
		\hline 
		\hline \\[-1.8ex] 
		\textit{Note:}  & \multicolumn{2}{r}{$^{*}$p$<$0.1; $^{**}$p$<$0.05; $^{***}$p$<$0.01} \\ 
	\end{tabular} 
\end{table} 

	\noindent Oil (negative): Holding all other variables constant, when the average ratio of fuel exports to total exports in 1984-86 increases to 50\% or above from less than 50\%, the log-odds of having a negative difference in GDP between year t and t-1 versus having no change in GDP increase by 4.784, though this effect is not statistically significant.
	
	\noindent Oil (positive): Holding all other variables constant, when the average ratio of fuel exports to total exports increases to 50\% or above from less than 50\%, the log-odds of having a positive difference in GDP between year t and t-1 versus having no change in GDP increase by 4.576, though this effect is not statistically significant.
	
	\noindent Democracy (negative): Holding all other variables constant, when a country moves from a non-democracy to a democracy, the log-odds of having a negative difference in GDP between year t and t-1 versus having no change in GDP increase by 1.379. This effect is is only significant when the alpha is 0.1.
	
	\noindent Democracy (positive): Holding all other variables constant, when a country moves from a non-democracy to a democracy, the log-odds of having a positive difference in GDP between year t and t-1 versus having no change in GDP increase by 1.769. This effect is significant at an alpha-level of 0.05.
	
	\noindent Constant (negative): When all variables equal 0, the baseline log-odds of having a negative difference in GDP between year t and t-1 versus having no change in GDP are 3.805.
	
	\noindent Constant (positive): When all variables are equal to 0, the baseline log-odds of having a positive difference in GDP between year t and t-1 versus having no change in GDP are equal to 4.534.
	
	\noindent Overall, taking into consideration which results are significant with an alpha of 0.05, it can be said that countries that are democracies have statistically significant higher log-odds of having a positive difference in GDP between year t and t-1 than of having no change in GDP. 
		
	\item Construct and interpret an ordered multinomial logit with \texttt{GDPWdiff} as the outcome variable, including the estimated cutoff points and coefficients.
	
	\noindent\textbf{My Answer:}
	
	\noindent To run the ordered multinomial logit, I removed the reference category and set the factored outcome variable to be ordered.
	
	\lstinputlisting[language=R, firstline=70, lastline=77]{PS03_RL.R}
	
\begin{table}[!htbp] \centering 
	\caption{Table of Coefficients: Ordered Multinomial Logit} 
	\label{} 
	\begin{tabular}{@{\extracolsep{5pt}}lc} 
		\\[-1.8ex]\hline 
		\hline \\[-1.8ex] 
		& \multicolumn{1}{c}{\textit{Dependent variable:}} \\ 
		\cline{2-2} 
		\\[-1.8ex] & GDPWdiff \\ 
		\hline \\[-1.8ex] 
		Oil & $-$0.199$^{*}$ \\ 
		& (0.116) \\ 
		& \\ 
		Democracy & 0.398$^{***}$ \\ 
		& (0.075) \\ 
		& \\ 
		\hline \\[-1.8ex] 
		Threshold 1|2 & -0.731 \\ 
		Threshold 2|3 & -0.71 \\ 
		Observations & 3,721 \\ 
		\hline 
		\hline \\[-1.8ex] 
		\textit{Note:}  & \multicolumn{1}{r}{$^{*}$p$<$0.1; $^{**}$p$<$0.05; $^{***}$p$<$0.01} \\ 
	\end{tabular} 
\end{table} 
	
	\noindent Oil: Holding democracy constant, when the average ratio of fuel exports to total exports in 1984-86 increases to 50\% or above, from less than 50\%, the log-odds of being in a higher category of difference in GDP between year t and t-1 decreases by 0.199. This effect is only significant with an alpha of 0.1.
	
	\noindent Democracy: Holding oil constant, when a country moves from being a non-democracy to a democracy, the log-odds of being in a higher category of difference in GDP between year t and t-1 increases by 0.398. This effect is significant at an alpha level of 0.01.
	
	\noindent Threshold (Cutpoints) 1-2: When the latent variable (propensity for GDP change) crosses -0.731, the predicted outcome of the difference in GDP between year t and t-1 moves from the "negative" category to "no change".
	
	\noindent Threshold (Cutpoints) 2-3: When the latent variable crosses -0.71, the predicted outcome of the difference in GDP between year t and t-1 moves from the "no change" category to "positive". 
	
	\noindent Overall, given these results, the threshold values do not appear to hold a great amount of substantive value. This is due to the fact that these values are incredibly close together, indicating that it is rare for the latent variable to fall between them in the "no change" category and that these ordinal categories might not be the best fit for the data.
	
		
\end{enumerate}

\section*{Question 2} 
\vspace{.25cm}

\noindent Consider the data set \texttt{MexicoMuniData.csv}, which includes municipal-level information from Mexico. The outcome of interest is the number of times the winning PAN presidential candidate in 2006 (\texttt{PAN.visits.06}) visited a district leading up to the 2009 federal elections, which is a count. Our main predictor of interest is whether the district was highly contested, or whether it was not (the PAN or their opponents have electoral security) in the previous federal elections during 2000 (\texttt{competitive.district}), which is binary (1=close/swing district, 0="safe seat"). We also include \texttt{marginality.06} (a measure of poverty) and \texttt{PAN.governor.06} (a dummy for whether the state has a PAN-affiliated governor) as additional control variables. 

\begin{enumerate}
	\item [(a)]
	Run a Poisson regression because the outcome is a count variable. Is there evidence that PAN presidential candidates visit swing districts more? Provide a test statistic and p-value.
	
	\lstinputlisting[language=R, firstline=97, lastline=101]{PS03_RL.R}
	
\begin{table}[!htbp] \centering 
	\caption{Table of Coefficients: Poisson Model} 
	\label{} 
	\begin{tabular}{@{\extracolsep{5pt}}lc} 
		\\[-1.8ex]\hline 
		\hline \\[-1.8ex] 
		& \multicolumn{1}{c}{\textit{Dependent variable:}} \\ 
		\cline{2-2} 
		\\[-1.8ex] & PAN Visits \\ 
		\hline \\[-1.8ex] 
		Competitive District & $-$0.081 \\ 
		& z = $-$0.477 \\ 
		& p = 0.634 \\ 
		& \\ 
		Poverty & $-$2.080 \\ 
		& z = $-$17.728 \\ 
		& p = 0.000$^{***}$ \\ 
		& \\ 
		PAN-Affilitated Governor & $-$0.312 \\ 
		& z = $-$1.869 \\ 
		& p = 0.062$^{*}$ \\ 
		& \\ 
		Constant & $-$3.810 \\ 
		& z = $-$17.156 \\ 
		& p = 0.000$^{***}$ \\ 
		& \\ 
		\hline \\[-1.8ex] 
		Observations & 2,407 \\ 
		Log Likelihood & $-$645.606 \\ 
		Akaike Inf. Crit. & 1,299.213 \\ 
		\hline 
		\hline \\[-1.8ex] 
		\textit{Note:}  & \multicolumn{1}{r}{$^{*}$p$<$0.1; $^{**}$p$<$0.05; $^{***}$p$<$0.01} \\ 
	\end{tabular} 
\end{table} 

	\noindent There is not evidence that PAN presidential candidates visit swing districts more. The z-statistic from the Wald's test was found to be -0.477 and the "competitive district" coefficient has a p-value of 0.634, indicating that the null hypothesis that being a swing district would not have an effect on the number of PAN presidential candidate visits cannot be rejected. The test statistic and p-value can also be interpreted as showing no evidence that the "competitive district" coefficient is significantly different from 0. 
	
	\item [(b)]
	Interpret the \texttt{marginality.06} and \texttt{PAN.governor.06} coefficients.
	
	\noindent Marginality.06: Holding all other variables constant, for every one-unit increase in the poverty measure (marginality.06) for a district, there is a decrease in the number of times the winning PAN presidental candidate visits that district by a multiplicative factor of e$^{-2.080}$ or 0.125. This effect is statistically significant with an alpha-level of 0.01.
	
	\noindent PAN.governor.06: Holding all other variables constant, when a district has a PAN-affiliated governor, the number of times the winning PAN presidential candidate visits that district decreases by a multiplicative factor of e$^{-0.312}$ or 0.732. This effect is only statistically significant with an alpha-level of 0.1.
	
	\noindent Given these results, as the poverty measure was found to be statistically significant with a p-value less than 0.05, it can be concluded that PAN presidential candidates visit districts higher on the poverty measure less often than those lower on the poverty scale. I will note that the directionality of this measure was not clarified and therefore we cannot interpret the results as the district being more or less impoverished, solely that as a district becomes higher in marginality.06 (poverty measure), the number of times the winning PAN presidential candidate visits that district decreases.
	
	\item [(c)]
	Provide the estimated mean number of visits from the winning PAN presidential candidate for a hypothetical district that was competitive (\texttt{competitive.district}=1), had an average poverty level (\texttt{marginality.06} = 0), and a PAN governor (\texttt{PAN.governor.06}=1).
	
	\lstinputlisting[language=R, firstline=115, lastline=128]{PS03_RL.R}
	
	\begin{verbatim}           
		0.01494818 0.01494827 
	\end{verbatim}
	
	\noindent This result can be interpreted as, when a district is competitive, has an average poverty level, and has a PAN governor, the estimated mean number of visits from the winning PAN presidential candidate is 0.0149. This number is not a count value, but due to the fact that it is the average of the number of visits it can be interpreted as a decimal value. 
	
\end{enumerate}

\end{document}
